\begin{thebibliography}{22}
\bibitem[1]{ref1} Dimitris Achlioptas. Database-friendly random projections: Johnson-lindenstrauss with binary coins. J. Comput. Syst. Sci., 66(4):671–687, 2003. doi:10.1016/S0022-0000(03)00025-4.
\bibitem[2]{ref2} Sagar Bisoyi, Krishnamoorthy Dinesh, and Jayalal Sarma. On pure space vs catalytic space. Theor. Comput. Sci., 921:112–126, 2022. doi:10.1016/j.tcs.2022.04.005.
\bibitem[3]{ref3} Harry Buhrman, Richard Cleve, Michal Koucký, Bruno Loff, and Florian Speelman. Computing with a full memory: catalytic space. In Symposium on Theory of Computing, STOC 2014, pages 857–866. ACM, 2014.
\bibitem[4]{ref4} Harry Buhrman, Michal Koucký, Bruno Loff, and Florian Speelman. Catalytic space: Nondeterminism and hierarchy. Theory Comput. Syst., 62(1):116–135, 2018. doi:10.1007/ s00224-017-9784-7.
\bibitem[5]{ref5} Jin-yi Cai, Venkatesan T. Chakaravarthy, and Dieter van Melkebeek. Time-space tradeoff in derandomizing probabilistic logspace. Theory Comput. Syst., 39(1):189–208, 2006. doi: 10.1007/S00224-005-1264-9.
\bibitem[6]{ref6} Kuan Cheng and William M. Hoza. Hitting sets give two-sided derandomization of small space. Theory of Computing, 18(21):1–32, 2022. doi:10.4086/toc.2022.v018a021.
\bibitem[7]{ref7} Gil Cohen, Dean Doron, Ori Sberlo, and Amnon Ta-Shma. Approximating iterated multiplication of stochastic matrices in small space. In Proceedings of the 55th Annual ACM Symposium on Theory of Computing, STOC 2023, pages 35–45. ACM, 2023.
\bibitem[8]{ref8} James Cook and Ian Mertz. Catalytic approaches to the tree evaluation problem. In Proccedings of the 52nd Annual ACM SIGACT Symposium on Theory of Computing, STOC 2020, pages 752–760. ACM, 2020.
\bibitem[9]{ref9} James Cook and Ian Mertz. Tree evaluation is in space o(log n · log log n). Electron. Colloquium Comput. Complex., TR23-174, 2023. arXiv:TR23-174.
\bibitem[10]{ref10} Samir Datta, Chetan Gupta, Rahul Jain, Vimal Raj Sharma, and Raghunath Tewari. Randomized and symmetric catalytic computation. In Computer Science – Theory and Applications – 15th International Computer Science Symposium in Russia, CSR 2020, pages 211–223. Springer, 2020.
\bibitem[11]{ref11} Dean Doron, Edward Pyne, and Roei Tell. Opening up the distinguisher: A hardness to randomness approach for BPL = L that uses properties of BPL. Electron. Colloquium Comput. Complex., TR23-208, 2023. arXiv:TR23-208.
\bibitem[12]{ref12} Chetan Gupta, Rahul Jain, Vimal Raj Sharma, and Raghunath Tewari. Unambiguous catalytic computation. In 39th IARCS Annual Conference on Foundations of Software Technology and Theoretical Computer Science, FSTTCS 2019, volume 150 of LIPIcs, pages 16:1–16:13, 2019.
\bibitem[13]{ref13} Daniel M. Kane, Raghu Meka, and Jelani Nelson. Almost optimal explicit johnson-lindenstrauss families. In Leslie Ann Goldberg, Klaus Jansen, R. Ravi, and José D. P. Rolim, editors, Approximation, Randomization, and Combinatorial Optimization. Algorithms and Techniques – 14th International Workshop, APPROX 2011, and 15th International Workshop, RANDOM 2011, Princeton, NJ, USA, August 17-19, 2011. Proceedings, volume 6845 of Lecture Notes in Computer Science, pages 628–639. Springer, 2011. doi:10.1007/978-3-642-22935-0\_53.
\bibitem[14]{ref14} Daniel M. Kane and Jelani Nelson. Sparser johnson-lindenstrauss transforms. J. ACM, 61(1):4:1–4:23, 2014. doi:10.1145/2559902.
\bibitem[15]{ref15} Ian Mertz. Reusing space: Techniques and open problems. Bulletin of EATCS, 141(3), 2023.
\bibitem[16]{ref16} Noam Nisan. Psuedorandom generators for space-bounded computation. In Harriet Ortiz, editor, Proceedings of the 22nd Annual ACM Symposium on Theory of Computing, May 13-17, 1990, Baltimore, Maryland, USA, pages 204–212. ACM, 1990. doi:10.1145/100216.100242.
\bibitem[17]{ref17} Noam Nisan. On read-once vs. multiple access to randomness in logspace. Theor. Comput. Sci., 107(1):135–144, 1993. doi:10.1016/0304-3975(93)90258-U.
\bibitem[18]{ref18} Noam Nisan. RL $\leq$ SC. Comput. Complex., 4:1–11, 1994. doi:10.1007/BF01205052.
\bibitem[19]{ref19} Aaron (Louie) Putterman and Edward Pyne. Near-optimal derandomization of medium-width branching programs. In Proceedings of the 55th Annual ACM Symposium on Theory of Computing, STOC 2023, pages 23–34, 2023.
\bibitem[20]{ref20} Edward Pyne, Ran Raz, and Wei Zhan. Certified hardness vs. randomness for log-space. In 64th IEEE Annual Symposium on Foundations of Computer Science, FOCS 2023, Santa Cruz, CA, USA, November 6-9, 2023, pages 989–1007. IEEE, 2023. doi:10.1109/FOCS57990.2023.00061.
\bibitem[21]{ref21} Michael E. Saks and Shiyu Zhou. $\mathsf{BP}_H\mathsf{Space}(S)\subseteq\mathsf{DSPACE}(S^{3/2})$. J. Comput. Syst. Sci., 58(2):376–403, 1999. doi:10.1006/jcss.1998.1616.
\bibitem[22]{ref22} D. Sivakumar. Algorithmic derandomization via complexity theory. In John H. Reif, editor, Proceedings on 34th Annual ACM Symposium on Theory of Computing, May 19-21, 2002, Montréal, Québec, Canada, pages 619–626. ACM, 2002. doi:10.1145/509907.509996.
\end{thebibliography}
